% Vietnamese translation for OI Public Library
% Source-PDF: IOI中国国家候选队论文/国家集训队2019论文集.pdf
\documentclass[11pt]{article}
\usepackage{oipl-vn}

\title{Tập luận văn đội tuyển quốc gia Trung Quốc năm 2019}
\author{Huấn luyện viên: Zhang Ruizhe}
\date{Tháng 5 năm 2019}

\begin{document}
\maketitle

\origtitle{Tập luận văn ứng viên đội tuyển quốc gia Trung Quốc Olympic Tin học năm 2019}
\sourcepath{IOI China candidate team papers / National training team 2019 collection.pdf}

\renewcommand{\abstractname}{Tóm tắt}
\begin{abstract}
PDF nguồn là tập hợp luận văn đội tuyển quốc gia Trung Quốc năm 2019. Tệp được
tạo bằng XeTeX và có lớp văn bản đọc được; tuy nhiên các trang có watermark
chéo, nên kết quả trích xuất xen thêm chữ từ watermark vào nội dung chính. Bản
LaTeX này chuyển ngữ phần nhận diện tập sách và mục lục, đồng thời ghi lại các
chủ đề chính để phục vụ bản dịch chi tiết hơn.
\end{abstract}

\section{Thông tin đầu sách}

Tập sách có tiêu đề \emph{Tập luận văn ứng viên đội tuyển quốc gia IOI Trung
Quốc năm 2019}. Trang bìa ghi huấn luyện viên là Zhang Ruizhe và thời điểm
phát hành là tháng 5 năm 2019. Tệp PDF gồm 231 trang.

\section{Mục lục đã dịch}

\begin{center}
\small
\begin{tabular}{r p{0.52\linewidth} p{0.24\linewidth}}
\textbf{Trang} & \textbf{Bài viết} & \textbf{Tác giả} \\
\hline
1 & Tính chất và ứng dụng của hai loại dãy truy hồi & Zhong Ziqian \\
17 & Bàn về bài toán bước đi ngẫu nhiên trên mô hình đồ thị & Wang Xiuhan \\
27 & Báo cáo ra đề \emph{Bài toán nhỏ dễ} & Yang Junzhao \\
38 & Bàn về bài toán tô màu đỉnh của đồ thị & Gao Jiaxuan \\
55 & Bàn về các bài toán liên quan tới đếm đường đi trên lưới & Dai Yan \\
69 & Mở rộng một số bài toán số học trong thi thuật toán và bước đầu về số nguyên Gauss & Li Jiaheng \\
89 & Báo cáo ra đề \emph{Bài luyện cơ bản về cây tròn-vuông} & Fan Zhiyuan \\
104 & Báo cáo ra đề \emph{Đếm điểm nguyên} và một số nghiên cứu về số nguyên Gauss & Xu Yixuan \\
122 & Bàn về thuật toán chia để trị trên cây & Zhang Zheyu \\
130 & Báo cáo ra đề \emph{Tổng tổ hợp} & Wu Siyang \\
141 & Bàn về một lớp cấu trúc dữ liệu ngắn gọn & Wang Siqi \\
152 & Thuật toán liên quan tới truy vấn chu kỳ xâu con và ứng dụng & Chen Sunli \\
165 & Báo cáo ra đề \emph{Công viên} & Wu Zuotong \\
192 & Bàn về cấu trúc dữ liệu có thể truy vết lịch sử & Kong Chaozhe \\
202 & Bàn về ứng dụng của ma trận Young trong thi tin học & Yuan Fangzhou
\end{tabular}
\end{center}

\section{Chủ đề chính}

Từ mục lục và phần đầu của bài thứ nhất, tập năm 2019 bao phủ các nhóm chủ đề
sau:
\begin{itemize}
  \item dãy truy hồi tuyến tính, dãy truy hồi đa thức và ứng dụng;
  \item bước đi ngẫu nhiên, tô màu đồ thị, chia để trị trên cây;
  \item đếm đường đi trên lưới, số nguyên Gauss và các bài toán số học;
  \item cây tròn-vuông, cấu trúc dữ liệu ngắn gọn và cấu trúc dữ liệu có lịch sử;
  \item truy vấn chu kỳ xâu con, ma trận Young và các báo cáo ra đề cụ thể.
\end{itemize}

\section{Ghi chú về trích xuất}

Phần thân của tập 2019 có thể trích xuất được mà không cần OCR. Bài đầu tiên
nêu rõ mục tiêu giới thiệu hai loại dãy truy hồi thường gặp: dãy truy hồi
tuyến tính và dãy truy hồi đa thức, cùng định nghĩa, tính chất, thuật toán và
ví dụ trong thi tin học. Khi mở rộng bản dịch, cần lọc bỏ watermark chéo trong
kết quả trích xuất để tránh lẫn chữ của bìa vào văn bản chính.

\end{document}
